\SetPractica{Ésteres}{Síntesis de ésteres mediante la esterificación de Fischer}

\section{Introducción teórica}



\section{Objetivos}

\begin{itemize}
    \item Sintetizar un éster a partir de un ácido carboxílico y un alcohol utilizando ácido sulfúrico como catalizador
    \item Comprobar la formación del éster mediante sus propiedades físicas y olor.
\end{itemize}

\section{Materiales, equipos y reactivos}

\begin{table}[H]
\centering{\small\setstretch{1.25}
\begin{tabular}{|m{0.3\linewidth}|m{0.3\linewidth}|m{0.3\linewidth}|} \hline
    
    \thead{Equipos} & \thead{Materiales} & \thead{Reactivos} \\ \hline
    
    {Termómetro} 
    
    &
    
    {Mechero Bunsen \par}
    {Soporte para mechero \par}
    {Malla de asbesto \par}
    {Matraz de fondo redondo \par}
    {Condensador de reflujo \par}
    {Probeta \par}
    {Embudo de separación \par}
    {Tubos de ensayo}
    
    & 
    
    {Agua destilada \par}
    {Ácido acético (\ce{CH3COOH}) \par}
    {Etanol \par}
    {Ácido sulfúrico (\ce{H2SO4}) \par}
    {Bicarbonato de sodio (\ce{NaHCO3})}
    
    \\ \hline
\end{tabular}}\end{table}

\section{Procedimientos}

\subsection{Preparación de la mezcla reactiva}
\begin{enumerate}
    \item En un matraz de fondo redondo de 100 mL coloca:
    \begin{enumerate}
        \item 10 mL de ácido acético
        \item 12 mL de etanol.
    \end{enumerate}
    
    \item Añade con cuidado unas gotas de ácido sulfúrico concentrado.
    \begin{enumerate}
        \item El ácido sulfúrico actuará como catalizador y ayudará a desplazar el equilibrio hacia la formación del éster.
    \end{enumerate}
    
\end{enumerate}

\subsection{Reflujo}
\begin{enumerate}
    \item Conecta el matraz a un condensador de reflujo
    \item Calienta suavemente la mezcla en un baño de agua a 60-70 °C durante 30-45 minutos.
    \begin{enumerate}
        \item El reflujo permite que los vapores del alcohol y el ácido se condensen y vuelvan a la mezcla sin perder reactivos.
    \end{enumerate}
    
    \item Mientras se calienta, asegúrate de que el sistema esté bien ventilado y que no haya pérdida de reactivos.
\end{enumerate}

\subsection{Enfriamiento de la mezcla}
\begin{enumerate}
    \item Después del tiempo de reflujo, retira el matraz del baño de agua
    \item Deja que se enfríe a temperatura ambiente.
\end{enumerate}

\subsection{Extracción del éster}
\begin{enumerate}
    \item Transfiere la mezcla a un embudo de separación
    \item Añade aproximadamente 20 mL de agua destilada.
    \item Agita suavemente el embudo de separación
    \item Deja reposar para permitir que se separen las fases.
    \begin{enumerate}
        \item El éster formará una capa orgánica, generalmente menos densa que el agua (dependiendo del éster específico, puede flotar o no).
    \end{enumerate}
\end{enumerate}

\subsection{Lavado de la fase orgánica}
\begin{enumerate}
    \item Lava la fase orgánica (que contiene el éster) con una solución de bicarbonato de sodio al 5\% para neutralizar cualquier ácido residual.
    \item Después de lavar con bicarbonato, lava nuevamente con agua destilada para eliminar cualquier resto de bicarbonato.
    \item Separa la fase acuosa (la inferior) de la fase orgánica (la superior, que contiene el éster).
\end{enumerate}

\subsection{Secado del éster}
\begin{enumerate}
    \item Añade una pequeña cantidad de un agente secante como sulfato de sodio anhidro a la fase orgánica y agita suavemente.
    \item Filtra para eliminar el agente secante y recupera el éster purificado.
\end{enumerate}

\subsection{Determinación de propiedades}
\begin{enumerate}
    \item Compara el olor del producto con los reactivos iniciales.
    \begin{enumerate}
        \item El éster formado (acetato de etilo, en este caso) tiene un olor afrutado característico.
    \end{enumerate}
\end{enumerate}

\subsection{Punto de ebullición}
\begin{enumerate}
    \item Determina el punto de ebullición del éster mediante destilación simple.
    \begin{enumerate}
        \item El acetato de etilo tiene un punto de ebullición de 77 °C.
    \end{enumerate}
\end{enumerate}